% Lab 6: Authentication, routing, and a live screen. Compile: latexmk -xelatex lab06.tex
\documentclass[10pt,aspectratio=169,t]{beamer}
\usetheme[lab]{upbminimal}

\course{Application Development for Mobile Devices}
\decklabel{Lab 6}
\title{Auth, routing \& a live screen}
\subtitle{Firebase Auth, go\_router guards, and a WebSocket screen}
\author{Dinu-Ștefan Rusu}
\institute{FILS, Universitatea Politehnica București}
\date{Week 12}

\begin{document}

\titleframe

\begin{frame}[eyebrow=Today]{What this lab is for}
  \vskip 2mm
  \begin{grid}[2]
    \cell[\faLock]{Gate the app on identity}{Wire Firebase Auth (email and password, plus Google) so authStateChanges drives what the user sees.}
    \cell[\faIcon{route}]{Guard routes, show live data}{Add go\_router with a redirect guard, then a screen that updates from a WebSocket without a refresh.}
  \end{grid}
  \vskip 6mm
  \taskmeta{Time}{2 hours}{Submit}{Branch \code{lab-6}, admd-labs repo, before next lab}
\end{frame}

\begin{frame}[fragile,eyebrow=Setup]{Before you write a line}
  \begin{steps}
    \item Branch from the Lab 5 state: \code{git checkout -b lab-6}
    \item Add packages: \code{firebase\_auth}, \code{google\_sign\_in}, \code{go\_router}, \code{web\_socket\_channel}.
    \item Run \code{flutterfire configure} again to regenerate \code{firebase\_options.dart}.
    \item Register your Android debug \emph{and} release SHA-1 fingerprints in the Firebase console. Google sign-in fails silently without them.
  \end{steps}
  \begin{hint}
\begin{shell}
keytool -list -v -keystore ~/.android/debug.keystore -alias androiddebugkey
\end{shell}
  \end{hint}
\end{frame}

\begin{frame}[fragile,eyebrow=Task I]{Sign-up and sign-in with Firebase Auth}
  \begin{columns}[T]
    \begin{column}{0.55\textwidth}
\begin{dart}
Future<void> signUp(String email, String pw) async {
  try {
    await FirebaseAuth.instance
        .createUserWithEmailAndPassword(
            email: email.trim(), password: pw);
  } on FirebaseAuthException catch (e) {
    throw mapAuthError(e.code);
  }
}
\end{dart}
    \end{column}
    \begin{column}{0.41\textwidth}
      \begin{steps}
        \item Email, password, a sign-up/sign-in toggle, and \code{mapAuthError} for one message per code.
      \end{steps}
    \end{column}
  \end{columns}
  \begin{takeaway}{One generic message for a failed sign-in.}
    Do not reveal whether an email is registered.
  \end{takeaway}
\end{frame}

\begin{frame}[fragile,eyebrow=Task I]{Google sign-in and the auth gate}
  \begin{columns}[T]
    \begin{column}{0.55\textwidth}
\begin{dart}
Future<UserCredential> signInWithGoogle() async {
  final account =
      await GoogleSignIn.instance.authenticate();
  final token = account.authentication.idToken;
  final cred = GoogleAuthProvider.credential(
      idToken: token);
  return FirebaseAuth.instance
      .signInWithCredential(cred);
}
\end{dart}
    \end{column}
    \begin{column}{0.41\textwidth}
      \kv{Done when}{}
      \begin{checklist}
        \item Google sign-in returns a signed-in user on a real device.
        \item Sign-out clears both sessions.
      \end{checklist}
      \vskip 3mm
      \muted{Call \code{initialize()} once at app start, before the first sign-in.}
    \end{column}
  \end{columns}
\end{frame}

\begin{frame}[fragile,eyebrow=Task I]{authStateChanges drives the app}
\begin{dart}
StreamBuilder<User?>(
  stream: FirebaseAuth.instance.authStateChanges(),
  builder: (context, snapshot) {
    if (snapshot.connectionState == ConnectionState.waiting) {
      return const SplashScreen();
    }
    final user = snapshot.data;
    return user == null ? const LoginScreen() : const HomeScreen();
  },
)
\end{dart}
  \begin{takeaway}{Listen to the stream. Never cache currentUser.}
    A cached field goes stale after a sign-out elsewhere or a token refresh. This stream is what Task II hooks the router into.
  \end{takeaway}
\end{frame}

\begin{frame}[fragile,eyebrow=Task II]{go\_router with a redirect guard}
  \begin{columns}[T]
    \begin{column}{0.6\textwidth}
\begin{dart}
final router = GoRouter(
  refreshListenable: GoRouterRefreshStream(
      FirebaseAuth.instance.authStateChanges()),
  redirect: (context, state) {
    final signedIn =
        FirebaseAuth.instance.currentUser != null;
    final atLogin = state.matchedLocation == '/login';
    if (!signedIn && !atLogin) return '/login';
    if (signedIn && atLogin) return '/';
    return null;
  },
);  // routes: [...]
\end{dart}
    \end{column}
    \begin{column}{0.36\textwidth}
      \begin{steps}
        \item Add \code{/login}, \code{/}, \code{/live} routes.
        \item Wire \code{refreshListenable} to \code{authStateChanges}.
        \item Sign-out button on Home calls \code{signOut()}.
      \end{steps}
    \end{column}
  \end{columns}
\end{frame}

\begin{frame}[eyebrow=Task II]{Why the guard belongs to the router}
  \begin{grid}[2]
    \cell[\faIcon{unlock}]{Without \code{refreshListenable}}{The redirect only runs on navigation. A background sign-out leaves the private screen on display until the user taps something.}
    \cell[\faIcon{sync}]{With it}{Every auth change re-runs \code{redirect} immediately, including the very first frame and any deep link.}
  \end{grid}
  \vskip 6mm
  \kv{Done when}{}
  \begin{checklist}
    \item Signing out from Home returns to \code{/login} without a manual navigation call.
    \item Opening \code{admd://live} while signed out lands on \code{/login}, not on \code{LiveScreen}.
    \item Signing in from \code{/login} lands back on \code{/}, never on \code{/login} itself.
  \end{checklist}
\end{frame}

\begin{frame}[fragile,eyebrow=Task III]{One screen, live messages}
\begin{dart}
final channel = WebSocketChannel.connect(
    Uri.parse('ws://10.0.2.2:8080/echo'));
await channel.ready;
StreamBuilder(
  stream: channel.stream,
  builder: (context, snapshot) => Text(
      snapshot.data as String? ?? 'Waiting...'),
)
\end{dart}
  \begin{takeaway}{await channel.ready is the line everyone forgets.}
    \code{connect()} returns instantly and never throws. Without \code{ready}, a rejected handshake surfaces later as a stream error.
  \end{takeaway}
\end{frame}

\begin{frame}[eyebrow=Task III]{Building the live screen}
  \begin{columns}[T]
    \begin{column}{0.55\textwidth}
      \begin{steps}
        \item Connect in \code{initState}, close in \code{dispose}.
        \item Keep the last N messages in a list; show them in a \code{ListView}.
        \item A text field and a send button that write to \code{channel.sink}.
        \item Show a connection-state banner: connecting, live, disconnected.
      \end{steps}
    \end{column}
    \begin{column}{0.41\textwidth}
      \kv{Done when}{}
      \begin{checklist}
        \item A message typed on one run appears without pressing refresh.
        \item Leaving the screen closes the socket (check the server log).
        \item The screen recovers after toggling airplane mode.
      \end{checklist}
    \end{column}
  \end{columns}
  \begin{hint}
    Android emulators reach your laptop's \code{localhost} at \code{10.0.2.2}, not \code{127.0.0.1}.
  \end{hint}
\end{frame}

\begin{frame}[fragile,eyebrow=Reference]{The echo server, in full}
\begin{golang}
var upgrader websocket.Upgrader
func main() {
  http.HandleFunc("/echo", func(w http.ResponseWriter, r *http.Request) {
    conn, _ := upgrader.Upgrade(w, r, nil)
    defer conn.Close()
    for {
      mt, msg, err := conn.ReadMessage()
      if err != nil { return }
      conn.WriteMessage(mt, msg)
    }
  })
  log.Fatal(http.ListenAndServe(":8080", nil))
}
\end{golang}
\end{frame}

\begin{frame}[fragile,eyebrow=Reference]{Why a local echo server}
  \begin{grid}[2]
    \cell[\faServer]{No public dependency}{A public echo endpoint can be slow, rate-limited, or offline during grading. This one runs on your laptop, on port 8080.}
    \cell[\faIcon{vial}]{You can see both sides}{The terminal running \code{go run .} prints every frame, which makes the WebSocket lifecycle from Lecture 11 concrete instead of abstract.}
  \end{grid}
  \vskip 4mm
\begin{shell}
go mod init echo
go get github.com/gorilla/websocket
go run .
\end{shell}
\end{frame}

\begin{frame}[eyebrow=Troubleshooting]{Errors and their fixes}
  \trouble{PlatformException(sign\_in\_failed, ApiException: 10)}{The SHA-1 fingerprint is missing or wrong in the Firebase console. Re-run \code{keytool} and re-download \code{google-services.json}.}
  \trouble{Too many redirects, or a blank screen at startup}{\code{redirect} returned the location it is already on. Return \code{null} to allow, never the current route.}
  \trouble{WebSocketChannelException: Connection refused}{The Go server is not running, or the emulator is using \code{127.0.0.1} instead of \code{10.0.2.2}.}
  \trouble{setState() called after dispose()}{A stream event arrived after the screen closed. Cancel the subscription in \code{dispose} before calling \code{sink.close}.}
\end{frame}

\closingframe{Push before you leave.}{Branch lab-6 · one commit per task · echo server included in the repo}

\end{document}
